% --------------------------------------
%  НАСТРОЙКИ МНОЖЕСТВЕННЫХ ФОРМ ДЛЯ РЕФЕРАТА
% --------------------------------------

\directlua{
  russian_plural = russian_plural or {}

  function russian_plural.form(number, one, few, many)
    local n = math.abs(tonumber(number) or 0)
    local last_two = n % 100
    local last = n % 10

    if last_two >= 11 and last_two <= 14 then
      return many
    end

    if last == 1 then
      return one
    end

    if last >= 2 and last <= 4 then
      return few
    end

    return many
  end
}

\newcommand{\russianPlural}[4]{%
  \directlua{
    tex.sprint(russian_plural.form(
      "\luaescapestring{#1}",
      "\luaescapestring{#2}",
      "\luaescapestring{#3}",
      "\luaescapestring{#4}"
    ))
  }%
}

\newcommand{\russianTables}[1]{%
  \russianPlural{#1}{таблицу}{таблицы}{таблиц}%
}

\newcommand{\russianAppendices}[1]{%
  \russianPlural{#1}{приложение}{приложения}{приложений}%
}

\newcommand{\russianSources}[1]{%
  \russianPlural{#1}{источник}{источника}{источников}%
}

\newcommand{\russianListings}[1]{%
  \russianPlural{#1}{листинг}{листинга}{листингов}%
}
